%% ============================================================================
%%  subdocumento-metodologias/contenido.tex
%%
%%  Ejemplo minimo de un subdocumento autocontenido. Sirve para verificar que
%%  el mecanismo funciona y como molde para los subdocumentos reales.
%%
%%  Reglas de la carpeta:
%%   - Este archivo SIEMPRE se llama contenido.tex y empieza por \chapter.
%%   - Las figuras van en figuras/ y se llaman por ruta relativa.
%%   - Los fragmentos van en tablas/ o donde quieras, y se traen con \parte{}.
%%   - NADA de aqui conoce la ruta de montaje: no escribas rutas que suban
%%     con ../ ni que empiecen desde la raiz del repositorio. Si lo haces,
%%     mover la carpeta rompe el capitulo.
%% ============================================================================
\setcounter{chapter}{5}
\chapter{Introducción a las metodologías del proyecto y software}
\label{cap:metodologias}

El capítulo presenta las metodologías de gestión de proyecto y de desarrollo de software que se adoptaran en el proyecto. La gestión del proyecto aplica la Guía del PMBOK, sexta edición (PMI, 2017), adaptada a un proyecto de 56 meses con dos etapas, solapamiento entre ellas y una operación que no puede detenerse; se usan los procesos de integración, alcance, cronograma, calidad, recursos, comunicaciones, riesgos, adquisiciones e interesados, con la profundidad que exige cada uno. El alcance y el cronograma tienen líneas base formales, porque los hitos del Formulario E-25 son fechas fijas y se aceptan por acta. El avance se mide con valor ganado y no con porcentajes declarados (RT-19.07). El flujo diario del equipo se gestiona con un tablero Kanban, que no reemplaza la línea base ni la aceptación formal de los entregables.

La sección 6.1 presenta el como LafroX gestionara el proyecto desde el inicio hasta la entrega final, ademas de gestionar los interesados, las comunicaciones y las adquisiciones. La sección 6.2 describe la metodología de desarrollo de software que se aplicara en el proyecto, incluyendo la planificación, el desarrollo iterativo, la gestión de cambios y la calidad del software.

Este capítulo

\section{Metodología de gestión de proyecto}

La gestión del proyecto sentara las bases para como el equipo se coordinara, planificara y ejecutara las entregas. Tener claro esos aspectos es vital para que el proyecto cumpla con los objetivos a tiempo siguiendo el cronograma contractual de 56 meses.

Gestionar el proyecto no basta solamente con planificar y coordinar el equipo, tambien consiste en la gestión de las personas que estan envueltas en el proyecto, la forma en la que nos comunicaremos con ellos y como manejaremos los contratos, ordenes de compra y acuerdos de nivel de servicio que se requieran para implementar la solución.

Para ello, nuestra gestión de proyecto se basa en la Guía del PMBOK, sexta edición (PMI, 2017), donde lo adaptamos a los 56 meses de contrato establecidos, con todas las 2 etapas, su solapamiento y la operación continua.

El proyecto se ejecuta mediante entregas incrementales dentro de las 2 etapas contractuales, priorizando la calidad y disponibilidad de la información.

Para limitar el trabajo en curso, se hara uso de tablero kanban, que muestra el trabajo pendiente, en ejecución, en revisión y terminado, delegando responsables a cada tarea, que tiene su definición de terminado y la evidencia a los criterios de aceptación.

\subsection{Gestión de interesados}

La gestión de interesados identifica a las personas, grupos y organizaciones que pueden influir en el proyecto o verse afectados por sus resultados. Para cada interesado se registran sus necesidades, expectativas, nivel de influencia, impacto potencial y disposición frente a los cambios. El análisis distingue quiénes participan en decisiones, quiénes aportan conocimiento operacional, quiénes validan entregables y quiénes son afectados o ejercen supervisión externa.

El jefe de proyecto mantiene actualizado el registro de interesados (paquete 1.4.1) y planifica su involucramiento de acuerdo con el alcance de cada entrega. La participación puede incluir levantamiento de necesidades, revisión de procesos, validación de criterios, pruebas y retroalimentación. Las inquietudes, resistencias, acuerdos y compromisos relevantes se registran con su responsable y seguimiento; los conflictos que no puedan resolverse en la instancia de trabajo se elevan a la autoridad correspondiente.

La tabla resume los principales grupos y roles identificados, sus intereses y la contribución esperada. La participación concreta se acuerda según las responsabilidades de cada interesado y las actividades de cada entrega; la inclusión en la tabla no implica que todos participen en todas las decisiones ni que tengan atribuciones de aprobación.

\subsection{Gestión de comunicaciones}

La gestión de las comunicaciones mantiene informados a los responsables de Puelche y a los demás interesados sobre los avances de las entregas que les competen, sostiene una comprensión común del estado del proyecto, facilita la toma de decisiones y comunica oportunamente las situaciones que puedan afectar los compromisos acordados.

Las reuniones convocan a las personas necesarias para analizar y resolver sus asuntos. Los acuerdos se registran en actas que identifican las decisiones, las actividades pendientes, sus responsables y los plazos; las actas de los comités se levantan dentro de los dos días hábiles siguientes (RT-19.09). La documentación vigente, los entregables, las actas y los registros de riesgos y de cambios se mantienen en el espacio colaborativo accesible al CLIENTE (RT-19.05).

Antes de cada entrega incremental se informa a los interesados pertinentes sobre los cambios previstos, las actividades de validación y el apoyo necesario para la puesta en funcionamiento. Las observaciones se registran y reciben respuesta, distinguiendo las que pueden atenderse dentro de lo acordado de las que requieren una solicitud de cambio.

\parte{tablas/comunicaciones}

\subsection{Gestión de adquisiciones}

La gestión de adquisiciones planifica, contrata y controla los bienes y servicios que requiere la solución: contratos, órdenes de compra, acuerdos con terceros y acuerdos de nivel de servicio. Cada adquisición tiene un responsable, una fecha de necesidad derivada del cronograma, una dependencia con los paquetes que la usan y una evidencia de recepción.

El CLIENTE compra el equipamiento de terreno y el de la sala técnica conforme a la especificación de cantidades y características que prepara LafroX (Caso 02, capítulo 11; SD3 E-09; paquetes 5.1.1 y 5.1.2). Esa compra no traslada al CLIENTE la instalación: LafroX habilita el recinto técnico en el espacio que proporciona el CLIENTE, instala y configura los equipos (fase 6 de la EDT) y contrata los servicios que le corresponden. Las licencias de terceros se constituyen a nombre del CLIENTE (5.2.3).

Cada fila tiene en el registro de adquisiciones su estado, su proveedor y su fecha comprometida. Un atraso que amenace el H3 o una ola se escala al Comité Ejecutivo con su análisis de impacto.

\subsection{Gestión de integración}

La gestión de integración mantiene alineados los componentes del proyecto para que las decisiones sobre alcance, cronograma, costos, recursos, calidad, riesgos y adquisiciones se analicen de manera coordinada. El jefe de proyecto consolida estos elementos en el plan para la dirección del proyecto, lo comunica a los responsables y lo actualiza cuando se aprueban cambios.

Durante la ejecución, el jefe de proyecto coordina las actividades y sus dependencias, revisa el avance de los entregables y gestiona los impedimentos que afectan a más de un equipo o área. Las decisiones, acuerdos, riesgos transversales y asuntos que requieran escalamiento quedan registrados con su responsable, resolución y fecha de seguimiento.

Toda solicitud de cambio que pueda afectar el alcance, los plazos, los costos o los criterios de aceptación se documenta y evalúa antes de implementarla. El análisis considera sus efectos en el plan, los entregables relacionados, los riesgos y la operación. El jefe de proyecto eleva la solicitud a la autoridad correspondiente de Puelche y comunica la decisión a las personas afectadas. Solo los cambios aprobados se incorporan a la planificación vigente. Las acciones correctivas destinadas a cumplir los compromisos aprobados se registran y supervisan; si modifican esos compromisos, se tramitan además como solicitudes de cambio.

Al finalizar cada entrega se comprueba que sus resultados satisfagan los criterios de aceptación y cuenten con la evidencia requerida. Las observaciones pendientes se registran con responsable y tratamiento acordado, y se documentan las lecciones aprendidas.

\subsection{Control del alcance, del cronograma y del valor ganado}

El alcance se controla contra la línea base y la matriz de trazabilidad del Formulario T-12: un requerimiento sólo se da por cumplido cuando pasa la prueba asignada en el paquete que lo construye. El cronograma se controla contra la red del Formulario T-15, sección 5, que incluye los diez días hábiles de revisión del CLIENTE antes de cada hito (Art. 18.3).

El avance se mide con valor ganado (RT-19.07). El valor planificado de cada mes es la suma de las horas hombre programadas de los paquetes en curso según el Formulario T-15; el valor ganado acredita las horas de un paquete sólo en hitos de avance definidos de antemano (inicio, revisión interna y aceptación) y no por porcentaje declarado; el costo real son las horas registradas. Con ellos se calculan el índice de desempeño del cronograma (SPI) y el del costo (CPI). Un SPI o un CPI bajo 0,90, o una desviación mayor al 10 %, obligan a presentar un plan de recuperación dentro de cinco días hábiles (paquete 1.8.6).

Las reservas se controlan por separado: la capacidad protegida de la Etapa 1 y las necesidades de contingencia del SD8 sólo se usan con autorización registrada, y su consumo se informa en el mismo informe mensual.

\subsection{Mecanismos de decisión y cadencias de gobierno}

Los comités del Art. 71° de las Bases Administrativas se organizan así:
\begin{itemize}
    \item Reunión de seguimiento semanal (paquete 1.3.5), con el jefe de proyecto y líderes de frente: revisa avance, bloqueos y datos para los comités; no reemplaza a ningún comité.
    \item Comité de Proyecto quincenal (paquete 1.8.3), con jefe de proyecto, líderes y Contraparte Técnica: revisa avance de entregables y el registro vivo de riesgos (RT-19.04).
    \item Comité Ejecutivo mensual (paquete 1.8.2), con gerencias del CLIENTE y dirección de LafroX: resuelve escalamientos, cambios de alcance o plazo y uso de reservas.
    \item Comité de Arquitectura mensual (paquete 1.8.4), con Arquitecto de Solución, Seguridad y TI del CLIENTE: decide asuntos de arquitectura y mantiene su registro.
    \item Comité de Operación mensual desde el mes 13 hasta el mes 56 (paquete 1.8.5), con Líder de Operación / SRE, mesa de servicio, Operaciones y TI del CLIENTE: revisa niveles de servicio, incidentes y problemas, capacidad de la mesa, pruebas de continuidad y plan de mejoras.
\end{itemize}

El Comité de Operación empieza con la primera marcha blanca, porque desde entonces la solución atiende operaciones reales. Revisa cada mes la disponibilidad de los servicios críticos, los incidentes críticos y altos y sus causas, la atención de la mesa (respuesta antes de 20 segundos, abandono y resolución al primer contacto), las pruebas de restauración y de continuidad y los cambios programados. Sus actas registran acuerdos, responsables y plazos (RT-19.09).




\section{Metodología de desarrollo de software}

La metodología de desarrollo se coordina con la gestión del proyecto y con sus entregas incrementales. Se utiliza \textbf{Rational Unified Process (RUP)}, un proceso iterativo e incremental que organiza el desarrollo en cuatro fases y permite gestionar requisitos, arquitectura, implementación y pruebas de manera progresiva. Su elección se basa en que el proyecto tiene requisitos regulatorios y de continuidad que exigen una arquitectura validada temprano, documentación trazable y entregas formales por etapa, sin renunciar a iteraciones cortas que reduzcan el tiempo hasta disponer de software verificable.

RUP se organiza en las siguientes fases:
\begin{enumerate}
    \item \textbf{Inicio}: se delimitan el alcance y los objetivos de la solución, se identifican los principales interesados y requisitos y se elaboran las estimaciones iniciales. También se reconocen los riesgos principales y se establece una visión inicial del producto.
    \item \textbf{Elaboración}: se profundizan los requisitos, incluidos los casos de uso, y se define la arquitectura de la solución. Se analizan los riesgos prioritarios y se prepara el plan de desarrollo, de modo que las decisiones de diseño más relevantes se validen antes de construir. Los prototipos con usuarios terminan antes de construir cada módulo.
    \item \textbf{Construcción}: se desarrollan las capacidades priorizadas mediante iteraciones. En cada iteración se analizan los requisitos correspondientes, se diseña e implementa la solución y se realizan pruebas. Los resultados se revisan con los interesados pertinentes y se ajustan cuando corresponde.
    \item \textbf{Transición}: se prepara la solución para su uso: pruebas de aceptación, resolución de observaciones, preparación de usuarios, marcha blanca y despliegue conforme a los hitos aprobados.
\end{enumerate}

Las iteraciones producen resultados verificables, pero no implican una puesta en producción o una marcha blanca. La marcha blanca y la aceptación formal se realizan en los momentos definidos para cada etapa contractual. Los casos de uso expresan y siguen los requisitos funcionales, y la matriz del Formulario T-12 los vincula con su prueba.

La secuencia de iteraciones se coordina con las dos etapas contractuales: en la Etapa 1, preventa, recepción, bodega, preparación y cross-docking; trazabilidad de lotes y de la cadena de frío; planificación de rutas; confirmación de recepción por los conductores; y manejo de efectivo y retiros. En la Etapa 2, los portales y el intercambio electrónico con los clientes y, luego, el costo de servir. La transición sigue el cronograma: marcha blanca de la Etapa 1 entre los meses 13 y 15 y paso a producción en el mes 16; marcha blanca de la Etapa 2 en los meses 19 y 20 y paso a producción en el mes 21.

\subsection{Arquitectura evolutiva, refactorización y deuda técnica}

La arquitectura se establece durante Elaboración y se aprueba en el H2; se valida progresivamente en las iteraciones de Construcción. Cuando los nuevos requisitos o los resultados de las pruebas justifican cambios, la arquitectura se ajusta mediante una nueva versión de la decisión en el registro de decisiones de arquitectura del SD4, con su fecha, estado y acta del Comité de Arquitectura.

Durante la construcción, el equipo refactoriza el software para mejorar su estructura interna y facilitar su mantenimiento, preservando su comportamiento funcional; las pruebas automatizadas comprueban que el comportamiento no cambió. Las limitaciones o compromisos técnicos que puedan dificultar cambios futuros se registran como deuda técnica, con su impacto y prioridad. Una deuda clasificada como bloqueante por el análisis estático impide el despliegue; las demás se planifican en las iteraciones y se revisan en el Comité de Arquitectura.

\subsection{DevSecOps, integración y entrega continuas, infraestructura como código y pruebas automatizadas}

Cada cambio pasa por el pipeline de integración continua que define el SD4, sección 4.2: GitLab CI orquesta y AWS CodeBuild construye cada imagen de forma hermética, con procedencia SLSA nivel 3. El pipeline ejecuta la auditoría de dependencias con \texttt{composer audit}, las pruebas PHPUnit, el análisis estático con PHPStan y Larastan, el formato con Laravel Pint, las pruebas de contrato contra OpenAPI 3.1 y AsyncAPI 2.6, el escaneo de secretos y de imágenes de contenedor y la medición de cobertura.

Las compuertas son bloqueantes, no opcionales. El pipeline detiene la promoción ante:
\begin{itemize}
    \item un hallazgo de seguridad crítico o alto en dependencias, código, secretos o imagen;
    \item un contrato público roto sin una nueva edición de la interfaz;
    \item una cobertura de la lógica de negocio inferior al 70\% (RT-04.11);
    \item una cobertura de líneas por pruebas unitarias del código modificado inferior al 80\%, política corporativa de LafroX (paquete 1.5.2);
    \item una deuda técnica bloqueante o una prueba en falla.
\end{itemize}

La imagen aprobada se firma, se publica en Elastic Container Registry y se promueve por su digest, de modo que ningún ambiente recompila. La infraestructura de los cinco ambientes se describe como código: Terraform administra los recursos con estados separados por ambiente y Ansible configura los hosts; cada recurso tiene un único propietario de código. Las migraciones de base de datos siguen la estrategia de expandir y contraer y declaran su reversión.

La entrega continua mantiene versiones verificadas y listas para su despliegue en cada ambiente. El paso a Producción se ejecuta sólo dentro de las ventanas permitidas por el caso y, durante el desarrollo, sólo en los hitos aprobados: ningún cambio se despliega en septiembre, en diciembre ni en los tres primeros días hábiles de un mes.

\subsection{Ceremonias, cadencias y decisiones del desarrollo}

El trabajo del equipo de desarrollo se organiza en iteraciones de Construcción de dos semanas. Al inicio de cada iteración se revisan los requisitos y prioridades, se seleccionan las tareas que caben en la capacidad registrada y se aclaran sus criterios de aceptación. Durante la iteración, el equipo sigue el avance en el tablero Kanban, con bloqueos, responsables y dependencias.

Al cierre de cada iteración se demuestra el incremento verificable y se contrasta con sus criterios de aceptación; cuando es pertinente, se presenta a los interesados para recoger observaciones. El equipo revisa además las dificultades y aprendizajes de la iteración y registra las acciones de mejora con su responsable. Ninguna demostración constituye por sí misma una entrega formal o una puesta en producción.

Los artefactos del desarrollo son los casos de uso, la matriz de trazabilidad del Formulario T-12, el registro de decisiones de arquitectura, los contratos de interfaz, el registro de deuda técnica y los informes de pruebas de cada iteración.

Las decisiones técnicas necesarias para implementar los requisitos se toman dentro del equipo responsable y se documentan cuando afectan la arquitectura, las integraciones, la seguridad o la operación; las que cambian la arquitectura pasan por el Comité de Arquitectura. Si una decisión modifica el alcance, el plazo, los costos o los criterios de aceptación aprobados, se gestiona como una solicitud de cambio según la sección 6.1.4.

\section*{Referencias}
\begin{itemize}
    \item Distribuidora Puelche S.A. (2026a). \textit{Bases Administrativas de Licitación N.\textsuperscript{o} TFEP-01/2026: Contratación de Solución Integral de Software y Servicios de Operación}.
    \item Distribuidora Puelche S.A. (2026b). \textit{Bases Técnicas Transversales de Licitación N.\textsuperscript{o} TFEP-01/2026}.
    \item Distribuidora Puelche S.A. (2026c). \textit{Caso 02: Logística. Especificaciones del problema y operación de Distribuidora Puelche S.A.}.
    \item Distribuidora Puelche S.A. (2026d). \textit{Aclaraciones de la Licitación N.\textsuperscript{o} TFEP-01/2026}.
    \item Kruchten, P. (2004). \textit{The Rational Unified Process: An introduction} (3.\textsuperscript{a} ed.). Addison-Wesley.
    \item Project Management Institute. (2017). \textit{La guía de los fundamentos para la dirección de proyectos (Guía del PMBOK\textsuperscript{\textregistered})} (6.\textsuperscript{a} ed.). Project Management Institute.
\end{itemize}

\section*{Declaración de uso de IA}

La tabla declara el uso de herramientas de inteligencia artificial en este subdocumento y en sus formularios; se consolida en el Formulario A-6.

\begin{tablalafrox}{Declaración de uso de inteligencia artificial}{tab:declaracion-uso-ia}{YYYYYY}
    {\cab{Sección} & \cab{Herramienta} & \cab{Finalidad del uso} & \cab{Nivel en texto} & \cab{Nivel en diagramas} & \cab{Revisión humana (quién y qué verificó)}}
Introducción & Claude Code & Redacción de la introducción y conexión con SD4 y SD7 (7 de octubre de 2026) & Alto & Ninguno & No documentada \\
6.1 Metodología de Gestión de Proyectos & Codex; Claude Code & Responsabilidades de adquisición y cadencias (6 de octubre de 2026); PMBOK adaptado, valor ganado, matriz de adquisiciones y Comité de Operación (7 de octubre de 2026) & Alto & Ninguno & No documentada \\
6.2 Metodología de Desarrollo Software & Claude Code & Compuertas DevSecOps, herramientas del SD4, cadencia de iteraciones y artefactos (7 de octubre de 2026) & Alto & Ninguno & No documentada \\
Formulario T-9 & Claude Code & Portada del formulario & Bajo & Ninguno & No documentada \\
Formulario T-10 & Claude Code & Portada del formulario & Bajo & Ninguno & No documentada \\
\end{tablalafrox}
